% fig_vl03_vollst.tex — Vollständigkeitsrelation: ein Vektor, zerlegt in
% zwei verschiedenen Orthonormalbasen desselben Raumes.
%
% Follows the conventions of fig_vl11_kasten.tex (the reference plate).
%
%   * geometry comes from the expansion derived in the lecture,
%         |Psi> = sum_j <phi_j|Psi> |phi_j>,
%     evaluated for the two orthonormal bases the lecture has already
%     introduced in C^2,
%         {|e_0>, |e_1>}   and   |e_±> = (|e_0> ± |e_1>)/sqrt(2),
%     on the concrete unit vector
%         |Psi> = (2|e_0> + |e_1>)/sqrt(5)   (real, so it fits in a plane).
%     Every coordinate below is that expansion, not a sketch:
%         c_0 = 2/sqrt5,  c_1 = 1/sqrt5   ->  |c|^2 = 4/5, 1/5
%         c_± = (2 ± 1)/sqrt10            ->  |c|^2 = 9/10, 1/10
%     and the two component rectangles therefore share |Psi> as their
%     diagonal exactly, because sum_j c_j |phi_j> = |Psi> in both bases.
%
%   * colour encodes *which basis*, and nothing else: cgblue for
%     {|e_0>,|e_1>}, birrenorange for {|e_+>,|e_->}.  The state |Psi>
%     itself is neutral navy — it belongs to no basis, which is the whole
%     point.  The coordinate rays are thin neutral grey: reference lines,
%     not data.
%
%   * within one basis the two members are told apart by tint (full
%     colour / 35%), not by a second hue, so the bars survive greyscale
%     printing; the numbers are printed as well, so nothing depends on
%     colour at all.  In (a) the same redundancy is carried by the dash
%     pattern — the first basis closes its rectangle with dashes, the
%     second with dots — and every arrow carries its own label, so the
%     two decompositions stay separable with the colour thrown away.
%     The neutral rays are held at birrengray!60: at full strength they
%     print at the same grey value as birrenorange and swallowed the
%     orange component vectors lying on top of them.
%
% Panel (a) the two decompositions of the same |Psi>.  Two different
%           rectangles, one diagonal.
% Panel (b) the weights |c_j|^2 those rectangles imply.  4/5 : 1/5 against
%           9/10 : 1/10 — the split is basis-dependent, the total is not.
%           This is the content of the completeness axiom
%           ||psi||^2 = sum_j |<phi_j|psi>|^2, and the reason the numbers
%           can later be read as probabilities.

\begin{center}
\begin{tikzpicture}[
    projvec/.style   ={-{Stealth[length=5pt,width=3.5pt]}, line width=1.3pt},
    statevec/.style  ={-{Stealth[length=6pt,width=4pt]}, line width=1.6pt,
                       spacecadet},
    basisray/.style  ={-{Stealth[length=4pt,width=3pt]}, line width=0.5pt,
                       birrengray!60},
    projdash/.style  ={line width=0.6pt, densely dashed},
    projdot/.style   ={line width=0.7pt, densely dotted},
]

% ─────────────────────────────────────────────────────────────
%  (a)  |Psi> resolved in two orthonormal bases
% ─────────────────────────────────────────────────────────────
\begin{axis}[
    name=geom,
    tpcplot,
    scale only axis,
    width=5.4cm, height=5.4cm,
    axis lines=none, grid=none,
    xmin=-0.38, xmax=1.38,
    ymin=-0.64, ymax=1.12,     % same span as x: 45 degrees looks like 45
    clip=false,
    title={(a) \tpcifenglish{One vector, two bases}%
                            {Ein Vektor, zwei Basen}},
]

  % ---- coordinate rays of the two bases (reference lines, neutral) ----
  \draw[basisray] (axis cs:-0.18,0)      -- (axis cs:1.24,0);
  \draw[basisray] (axis cs:0,-0.18)      -- (axis cs:0,1.00);
  \draw[basisray] (axis cs:-0.127,-0.127) -- (axis cs:0.721,0.721);
  \draw[basisray] (axis cs:-0.127,0.127)  -- (axis cs:0.424,-0.424);

  \node[lbl, anchor=west]       at (axis cs:1.27,0)      {$|e_0\rangle$};
  \node[lbl, anchor=south]      at (axis cs:0,1.03)      {$|e_1\rangle$};
  \node[lbl, anchor=south west] at (axis cs:0.74,0.74)   {$|e_+\rangle$};
  \node[lbl, anchor=north west] at (axis cs:0.44,-0.44)  {$|e_-\rangle$};

  % ---- the two component rectangles: same diagonal, different sides ----
  % basis {e_0,e_1}:  feet at (2/sqrt5, 0) and (0, 1/sqrt5)
  \draw[projdash, cgblue] (axis cs:0.894427,0)
                       -- (axis cs:0.894427,0.447214);
  \draw[projdash, cgblue] (axis cs:0,0.447214)
                       -- (axis cs:0.894427,0.447214);
  % basis {e_+,e_-}:  feet at 3/sqrt10 along e_+ and 1/sqrt10 along e_-
  \draw[projdot, birrenorange] (axis cs:0.670820,0.670820)
                            -- (axis cs:0.894427,0.447214);
  \draw[projdot, birrenorange] (axis cs:0.223607,-0.223607)
                            -- (axis cs:0.894427,0.447214);

  % ---- the component vectors c_j |phi_j> ----
  \draw[projvec, cgblue]       (axis cs:0,0) -- (axis cs:0.894427,0);
  \draw[projvec, cgblue]       (axis cs:0,0) -- (axis cs:0,0.447214);
  \draw[projvec, birrenorange] (axis cs:0,0) -- (axis cs:0.670820,0.670820);
  \draw[projvec, birrenorange] (axis cs:0,0) -- (axis cs:0.223607,-0.223607);

  % ---- the state itself ----
  \draw[statevec] (axis cs:0,0) -- (axis cs:0.894427,0.447214);
  \node[lbl, text=spacecadet, font=\small, anchor=south west]
    at (axis cs:0.91,0.47) {$|\Psi\rangle$};

  % ---- component labels, each beside its own arrow ----
  \node[lbl, text=cgblue, font=\scriptsize, anchor=north]
    at (axis cs:0.72,-0.045) {$c_0|e_0\rangle$};
  \node[lbl, text=cgblue, font=\scriptsize, anchor=east]
    at (axis cs:-0.04,0.22)  {$c_1|e_1\rangle$};
  \node[lbl, text=birrenorange, font=\scriptsize, anchor=south east]
    at (axis cs:0.58,0.685)  {$c_+|e_+\rangle$};
  \node[lbl, text=birrenorange, font=\scriptsize, anchor=north east]
    at (axis cs:0.155,-0.255) {$c_-|e_-\rangle$};

  \node[lbl, text=birrengray, font=\scriptsize, anchor=north east]
    at (axis cs:1.38,1.12) {$c_j \equiv \langle \phi_j|\Psi\rangle$};
\end{axis}

% ─────────────────────────────────────────────────────────────
%  (b)  The weights those two decompositions carry
% ─────────────────────────────────────────────────────────────
\begin{axis}[
    tpcplot,
    scale only axis,
    at={(geom.north east)}, anchor=north west, xshift=22mm,
    width=5.0cm, height=4.2cm,
    axis lines=left, grid=none,
    xmin=0, xmax=1.58, ymin=0.42, ymax=2.55,
    xtick={0,0.5,1}, xticklabels={$0$,$1/2$,$1$},
    ytick={1,2},
    yticklabels={{$\{|e_+\rangle,|e_-\rangle\}$},{$\{|e_0\rangle,|e_1\rangle\}$}},
    ytick style={draw=none},
    xlabel={$|c_j|^2$},
    clip=false,
    title={(b) \tpcifenglish{Weights}{Gewichte}},
]

  % the common total, drawn first so it does not sit on the bar borders
  \draw[levellinegrey] (axis cs:1,0.50) -- (axis cs:1,2.46);

  % basis {e_0,e_1}:  4/5 and 1/5
  \fill[cgblue]     (axis cs:0,1.74)  rectangle (axis cs:0.8,2.26);
  \fill[cgblue!35]  (axis cs:0.8,1.74) rectangle (axis cs:1.0,2.26);
  \draw[spacecadet, line width=0.5pt]
                    (axis cs:0,1.74)  rectangle (axis cs:1.0,2.26);

  % basis {e_+,e_-}:  9/10 and 1/10
  \fill[birrenorange]    (axis cs:0,0.74)  rectangle (axis cs:0.9,1.26);
  \fill[birrenorange!35] (axis cs:0.9,0.74) rectangle (axis cs:1.0,1.26);
  \draw[spacecadet, line width=0.5pt]
                         (axis cs:0,0.74) rectangle (axis cs:1.0,1.26);

  \node[font=\footnotesize, text=white] at (axis cs:0.40,2.00)
    {$|c_0|^2=\tfrac45$};
  \node[lbl, text=cgblue, font=\footnotesize, anchor=west]
    at (axis cs:1.04,2.00) {$|c_1|^2=\tfrac15$};

  \node[font=\footnotesize, text=white] at (axis cs:0.45,1.00)
    {$|c_+|^2=\tfrac{9}{10}$};
  \node[lbl, text=birrenorange, font=\footnotesize, anchor=west]
    at (axis cs:1.04,1.00) {$|c_-|^2=\tfrac{1}{10}$};

  \node[lbl, text=birrengray, font=\scriptsize]
    at (axis cs:0.50,1.50)
    {$\textstyle\sum_j |c_j|^2 = \|\Psi\|^2 = 1$};
\end{axis}

\end{tikzpicture}
\end{center}
